\documentclass[article]{memoir}
\usepackage{polyglossia}[2022/10/26]
\setdefaultlanguage[forceheadingpunctuation]{hungarian}

\usepackage{csquotes, amsmath, amsthm, fixme, graphicx, datatool, nevelok}
\swapnumbers %% a babel megfordítja. A polyglossia nem. De a babel pontot is tesz a számcimke után
\theoremstyle{plain}
\newtheorem{proposition}{állítás}[chapter]
\newcommand{\ip}[2]{\langle#1,#2\rangle}
\begin{document}
Ide másoltam a könyvből a Bessel--Schwartz páros igazolását.
A Schwartzot a Besselből kapjuk azt egyetlen vektorra alkalmazva.
Ha ez így van, akkor persze ha felírjuk a Bessel igazolását egyetlen vektorra, akkor is a Schwartz-ot kell kapnunk.
Talán ilyesmi lehet a Schwartz-egyenlőtlenség rövid igazolásának hátterében.

\begin{proposition}[Bessel]\index{Bessel-egyenlőtlenség}\label{pr:Bessel}
	Legyen $\left( V,\ip{\cdot}{\cdot} \right)$ skalárisszorzatos-tér.
	Tegyük fel, hogy $\left\{ x_1,\ldots,x_k \right\}$ egy ortonormált rendszer.
	Ekkor minden $u\in V$ vektorra:
	\begin{enumerate}
		\item
		      fennáll a  \emph{Bessel-egyenlőtlenség:}
		      $\sum_{i=1}^k|\ip{u}{x_i}|^2\leq\ip{u}{u}$;
		\item
		      Az
		      \(
		      \hat{u}=u-\sum_{i=1}^k\ip{u}{x_i}x_i
		      \)
		      vektorra $\hat{u}\perp x_j$
		      ($j=1,\ldots,k$), azaz
		      \(
		      \hat{u}\in
		      \left\{ x_1,\ldots,x_k \right\}^\perp
		      %=
		      %\left(\lin\left\{ x_1,\ldots,x_k \right\} \right)^\perp
		      .
		      \)
		      \qedhere
	\end{enumerate}
\end{proposition}
\begin{proof}
	Jelölje a bizonyítás erejéig
	\begin{math}
		\alpha_j=\ip{u}{x_j}.
	\end{math}
	Ekkor az $\hat{u}=u-\sum_{j=1}^k\alpha_jx_j$ vektorra
	\begin{multline*}
		0
		\leq
		\ip{\hat{u}}{\hat{u}}
		=
		\ip{u}{u}+\ip{u}{-\sum_{j=1}^k\alpha_jx_j}+\ip{-\sum_{j=1}^k\alpha_jx_j}{u}+\ip{-\sum_{j=1}^k\alpha_jx_j}{-\sum_{i=1}^k\alpha_ix_i}
		\\
		=
		\ip{u}{u}-\sum_{j=1}^k\overline{\alpha_j}\ip{u}{x_j}-\sum_{j=1}^k\alpha_j\ip{x_j}{u}+\sum_{j=1}^k\sum_{i=1}^k\alpha_j\overline{\alpha_i}\ip{x_j}{x_i}
		\\
		=
		\ip{u}{u}-\sum_{j=1}^k\overline{\alpha_j}\alpha_j-\sum_{j=1}^k\alpha_j\overline{\alpha_j}+\sum_{j=1}^k\alpha_j\overline{\alpha_j}
		=\ip{u}{u}-\sum_{j=1}^k|\alpha_j|^2.
	\end{multline*}
	Ez éppen a Bessel-egyenlőtlenség.
	A merőlegességre vonatkozó második állítás már sokkal egyszerűbb:
	\[
		\ip{\hat{u}}{x_j}
		=
		\ip{u-\sum_{i=1}^k\alpha_ix_i}{x_j}
		=
		\ip{u}{x_j}-\sum_{i=1}^k\alpha_i\ip{x_i}{x_j}
		=
		\alpha_j-\alpha_j
		=
		0.
		\qedhere
	\]
\end{proof}
\begin{proposition}[Schwartz-egyenlőtlenség]\index{Schwartz-egyenlőtlenség}\label{pr:Schwartz-egyenlőtlenség}
	Legyen $\left( V,\ip{\cdot}{\cdot} \right)$ skalárisszorzatos-tér,
	$u,v\in V$.
	Ekkor
	\[
		|\ip{u}{v}|^2\leq \ip{u}{u}\ip{v}{v}.\qedhere
	\]
\end{proposition}
\begin{proof}
	Ha $v=0$, akkor mindkét oldal zérus.
	Ha $v\neq 0$, akkor a $\left\{ \frac{1}{\sqrt{\ip{v}{v}}}v \right\}$ rendszer egyetlen elemű ortonormált rendszer.
	Alkalmazva a Bessel-egyen\-lőt\-len\-sé\-get\index{Bessel-egyenlőtlenség} $k=1$ mellett
	\[
		\frac{1}{\ip{v}{v}}
		\left|\ip{u}{v}\right|^2
		=
		|\frac{1}{\sqrt{\ip{v}{v}}}
		\ip{u}{v}|^2
		=
		|\ip{u}{\frac{1}{\sqrt{\ip{v}{v}}}v}|^2
		\leq \ip{u}{u},
	\]
	ami kishíján a bizonyítandó állítás.
\end{proof}

Eddig volt a könyvbeli rész. 
Nézzük, akkor az alternatív bizonyítást, ami egyszerűen a Bessel-beli számolás egyetlen vektor mellett.
\begin{proof}
Tegyük fel először, hogy $\ip{v}{v}=1$. 
Jelölje $\alpha=\ip{u}{v}$ és $\hat{u}=u-\alpha v$.
Ekkor
	\begin{multline*}
		0
		\leq
		\ip{\hat{u}}{\hat{u}}
		=
		\ip{u}{u}+\ip{u}{-\alpha v}+\ip{-\alpha v}{u}+\ip{-\alpha v}{-\alpha v}
		\\
		=
		\ip{u}{u}-\overline{\alpha}\ip{u}{v}-\alpha\ip{v}{u}+\alpha\overline{\alpha}\ip{v}{v}
		\\
		=
		\ip{u}{u}-\overline{\alpha}\alpha-\alpha\overline{\alpha}+\alpha\overline{\alpha}
		=\ip{u}{u}-|\alpha|^2,
	\end{multline*}
    ami a bizonyítandó állítás a $\ip{v}{v}=1$ speciális esetben.
    Nem zérus $v$ mellett általános esetet kapjuk, ha felírjuk a már igazolt állítás a $\frac{1}{\sqrt{\ip{v}{v}}}v$ vektorra.
\end{proof}
\end{document}
